\documentclass[11pt]{ctexart}
\usepackage[a4paper,margin=1in]{geometry}
\usepackage{graphicx}
\usepackage{xcolor}
\usepackage{hyperref}
\definecolor{refgray}{gray}{0.45}
\title{\textbf{Market Views｜全球投行叙事汇编}\\[0.4em]{\large AI生产率通胀影响分歧加剧，存储芯片定价分化与消费量价背离成为今日跨主题主线}}
\author{KC桌面 自动生成}
\date{260729}
\begin{document}
\maketitle
\tableofcontents
\newpage
\section{一页摘要}
\begin{itemize}
  \item AI生产率对通胀的影响方向存在根本分歧：NOM认为需求渠道可能推高通胀，而美联储声明隐含生产率将结构性压低价格。
  \item 存储芯片定价分化显著：DDR4三季度月涨幅可超20\%，而主流DDR5涨价趋缓，市场关注点从整体周期转向结构性机会。
  \item 全球消费呈现量价背离：可口可乐销量增长5\%为17年最强，但中国高端消费（老铺黄金）受金价波动影响，净利润低于预期30\%。
  \item 中国汽车行业分化加速：成本压力在三季度集中显现，垂直整合能力强的车企与依赖外采的车企利润表现将显著分化。
  \item 代币化资产预计2030年达5.5万亿美元，但转型路径（混合模式 vs. 完全链上化）存在分歧，短期主导权未定。
  \item 依沃西单抗全球注册路径时间窗口拉长，关键OS数据读出节点在1H27，但Summit现金消耗速度要求提前补充资金。
  \item 全球人口总量拐点可能早于联合国预测的2084年，在2050年代前后到来，对消费需求的影响时点存在分歧。
  \item 中国工业利润高度集中于AI与上游原材料，下游消费持续承压，利润增速对AI资本开支和大宗商品价格敏感度高。
\end{itemize}
\section{全球投行叙事汇编}
今日纳入 64 篇投行/券商研究，按 10 个市场、资产与产业主题系统整理。
\newpage
\subsection{宏观与利率：AI生产率、汇率与政策不确定性}
\textbf{市场对AI生产率的通胀影响存在分歧，美联储宽松倾向与汇率管理、财政不确定性共同构成当前宏观主线。}
\subsubsection*{主流共识}
\begin{itemize}
  \item AI生产率提升对通胀的影响方向不确定，历史相关性不稳定，当前资本开支已推高设备价格。
  \item 美联储主席沃什对生产率的乐观立场使其保持宽松倾向，而非推动预防性加息。
  \item 韩国市场去杠杆进程接近完成，杠杆ETF强制平仓已结束，对冲基金去杠杆约90\%。
\end{itemize}
\subsubsection*{分歧与条件差异}
\begin{itemize}
  \item AI生产率是否必然带来通胀下降：NOM认为需求渠道可能推高通胀，而美联储声明隐含生产率将结构性压低价格。
  \item 法国预算审议结果：MS基准预测预算通过但赤字维持5.2\%，但若失败赤字可能升至5.9\%，市场定价尚未充分反映。
  \item 中国离岸信托新规对香港楼市影响：JPM认为影响不足2\%，但若扩大至所有内地税务居民则影响面更广。
\end{itemize}
\subsubsection*{机构观点}
\paragraph{NOM} AI生产率对通胀的影响方向不确定，历史相关性不稳定。当前AI资本开支已推高设备价格，超大规模云服务商资本开支持续扩张，消费电子价格出现上涨趋势。美联储主席沃什对生产率的乐观立场更可能使其保持宽松倾向，而非推动预防性加息。设备价格快速上涨时，企业有强烈动机提前下单锁定供应，形成需求端正向冲击。
{\small\color{refgray}数据：AI设备价格快速上涨，消费电子领域出现价格不确定性\par}
{\small\color{refgray}数据：超大规模云服务商资本开支计划持续扩张\par}
{\small\color{refgray}数据：设备价格快速上涨使债务融资实际利率深度为负\par}
{\small\color{refgray}边际变化：从传统生产率降低通胀的叙事转向强调需求渠道和资本开支对通胀的推升作用。\par}
\paragraph{NOM} USD/CNY中间价模型预测7月29日定价6.7640，较前一日模型预测下降288个基点。逆周期因子调整后为6.7727，较无因子版本高出87个基点，起到缓冲作用。隔夜美元指数变动、离岸人民币汇率波动是模型预测下移的主要驱动因素。逆周期因子差值扩大表明官方主动管理预期。
{\small\color{refgray}数据：模型预测6.7640，较前一日6.7928下降288个基点\par}
{\small\color{refgray}数据：含逆周期因子调整值6.7727，较前一日实际中间价下降201个基点\par}
{\small\color{refgray}数据：逆周期因子提供87个基点的缓冲\par}
{\small\color{refgray}边际变化：模型预测大幅下移，逆周期因子持续平滑汇率波动节奏。\par}
\paragraph{JPM} 中国离岸信托个税新规（20\%税率）对香港住宅市场方向性谨慎，但实际影响有限。通过内地控制的离岸信托购买的香港住宅交易占比不足2\%，不足以改变当前房价上升周期。新规追溯至2023年，反避税条款覆盖已获外国国籍者。启德新区内地买家比例较高，相对承压。
{\small\color{refgray}数据：公司买家占香港住宅交易4.5\%，其中内地控制不到一半\par}
{\small\color{refgray}数据：极端假设下销售减少不足2\%\par}
{\small\color{refgray}数据：租金收入新规仅补差额约5\%\par}
{\small\color{refgray}边际变化：新规方向性谨慎但影响可控，市场此前高估了冲击幅度。\par}
\paragraph{JPM} 外汇波动重新成为欧洲食品零售板块关键驱动因素。Jeronimo Martins 70\%收入来自波兰兹罗提，Ahold Delhaize 60\%来自美元，家乐福20\%来自巴西雷亚尔。家乐福股价与巴西雷亚尔/欧元高度相关，反映巴西业务贡献33\% EBIT。但汇率驱动的机械性盈利增长不应与基本面改善混为一谈。
{\small\color{refgray}数据：Jeronimo Martins 85\% EBIT来自外币业务\par}
{\small\color{refgray}数据：Ahold Delhaize 65\% EBIT来自外币业务\par}
{\small\color{refgray}数据：家乐福巴西业务贡献约33\% EBIT\par}
{\small\color{refgray}边际变化：首次系统梳理欧洲食品零售企业外汇敞口结构，区分折算影响与宏观环境影响。\par}
\paragraph{JPM} 韩国KOSPI自6月22日高点下跌约40\%，杠杆ETF强制平仓已完成，对冲基金去杠杆约90\%。杠杆ETF资产规模从500亿美元降至170亿美元，对冲基金多空比率从5.7倍降至3.2倍。KOSPI远期市盈率降至5倍，自由现金流回报处于危机水平。外资流出速度放缓，仓位结构从极端集中向分散回归。
{\small\color{refgray}数据：KOSPI自6月22日高点下跌约40\%\par}
{\small\color{refgray}数据：杠杆ETF资产规模从500亿美元降至170亿美元\par}
{\small\color{refgray}数据：对冲基金多空比率从5.7倍降至3.2倍\par}
{\small\color{refgray}数据：KOSPI远期市盈率降至5倍\par}
{\small\color{refgray}边际变化：去杠杆进程接近尾声，最紧张阶段可能已过。\par}
\paragraph{JPM} 苏泊尔二季度收入同比下滑3\%，低于预期增长4\%。出口是主要拖累，母公司SEB控制库存可能延续至下半年。出口毛利率大幅下滑，成本锁定未能对冲石化成本上涨。JPM将2026-2027年每股盈利预测下调5-9\%，下半年经营利润预计同比下降3\%。国内业务稳健但不足以完全对冲出口端压力。
{\small\color{refgray}数据：二季度收入同比下滑3\%，低于预期增长4\%\par}
{\small\color{refgray}数据：二季度经营利润同比下降16\%\par}
{\small\color{refgray}数据：出口毛利率假设从19\%下调至16.3\%\par}
{\small\color{refgray}数据：2026年每股分红预测从2.51元调整至2.01元\par}
{\small\color{refgray}边际变化：从盈利稳定溢价转向短期业绩不确定性，评级从超配下调至中性。\par}
\paragraph{MS} 法国未来12个月将经历预算审议、总统选举和可能提前议会选举三重事件叠加。2027年预算审议是市场定价第一个时间节点，基准预测预算通过但赤字维持GDP的5.2\%，若失败赤字可能升至5.9\%。总统选举结果在2027年初之前难以预判，历史显示民调与实际结果差距较大。政策不确定性将持续高企，但部分已在债券和股票中计入。
{\small\color{refgray}数据：基准预测赤字维持GDP的5.2\%\par}
{\small\color{refgray}数据：若预算失败赤字可能升至5.9\%\par}
{\small\color{refgray}数据：过去六次总统选举第一轮平均得票率25\%\par}
{\small\color{refgray}边际变化：从单一选举风险转向预算-选举-议会三重事件叠加框架。\par}
\subsubsection*{关键数据}
\begin{itemize}
  \item AI设备价格快速上涨，消费电子领域出现价格不确定性
  \item USD/CNY中间价模型预测6.7640，较前一日下降288个基点
  \item 香港住宅交易中内地离岸信托买家占比不足2\%
  \item 家乐福巴西业务贡献约33\% EBIT
  \item KOSPI远期市盈率降至5倍，杠杆ETF资产规模从500亿降至170亿美元
  \item 苏泊尔二季度经营利润同比下降16\%
  \item 法国基准赤字预测GDP的5.2\%，若预算失败升至5.9\%
\end{itemize}
\begin{figure}[htbp]
\centering
\includegraphics[width=0.88\linewidth]{figures/fig_065.jpg}
\caption{Figure 2 - Figure 1: HK residential property transactions – \% of company buyers \#\# "Mainland Chinese buyers" in HK residential market}
\end{figure}
\begin{figure}[htbp]
\centering
\includegraphics[width=0.88\linewidth]{figures/fig_078.jpg}
\caption{Figure 2 - Figure 3: Carrefour share price vs BRLEUR Currency – 5Y Chart Figure 4: Ahold share price vs USDEUR Currency – 5Y Chart \# Stock-specific exposure}
\end{figure}
\begin{figure}[htbp]
\centering
\includegraphics[width=0.88\linewidth]{figures/fig_089.jpg}
\caption{Figure 1 - Figure 1: KOSPI index has now declined \textbackslash{}\textasciitilde{}40\% from the peak of 22 Jun Figure 2: This week's move has brought index RSI to oversold levels Figure 3: KOSPI fwd P/E at 5x is cheap – even after adjusting for cyclicality}
\end{figure}
\begin{figure}[htbp]
\centering
\includegraphics[width=0.88\linewidth]{figures/fig_097.jpg}
\caption{Exhibit 2 - Exhibit 3: Since 1995, on average there have been 12 candidates in the first round \#\# Polls}
\end{figure}
{\small\color{refgray}本节综合 7 篇研究；涉及 NOM、JPM、MS。}
\newpage
\subsection{AI基础设施与应用：从芯片到系统，从替代到增强}
\textbf{AI投资正从单芯片性能竞争转向系统级方案与基础设施效率的比拼，同时AI对就业的影响呈现‘增强多于替代’的结构性分化，网络安全预算因AI驱动加速增长。}
\subsubsection*{主流共识}
\begin{itemize}
  \item AI基础设施投资重心从芯片规格转向系统级互联与光通信，中国厂商在超节点互联技术取得进展，华为、壁仞等推出大规模加速器互联方案。
  \item AI对就业市场的影响以增强为主，替代风险集中在常规文职岗位，印度非农就业中42-48\%可能被AI增强，仅8-12\%面临替代不确定性。
  \item 网络安全预算增速持续跑赢整体IT支出，AI成为核心驱动因素，身份安全（特别是非人类身份管理）和安全运营自动化成为新焦点。
\end{itemize}
\subsubsection*{分歧与条件差异}
\begin{itemize}
  \item PC市场出现‘量缩价升’结构性背离：出货量同比下滑5\%，但收入因高端化与涨价仍在增长，下半年出货量或进一步下滑。
  \item AI推理需求推动通用服务器CPU/GPU部署比例从1:3-4升至接近1:1，但供应瓶颈（基板、芯片）比预期更严重，部分厂商面临历史最严重供应限制。
  \item 光通信需求当前仅反映横向扩展，纵向扩展与CPO尚未贡献收入，但若落地将带来额外上行空间，时间节点不确定。
\end{itemize}
\subsubsection*{机构观点}
\paragraph{BofA} 康宁企业光通信收入同比增长65\%，AI贡献弹性高，但Q3指引略低于预期。当前估值超60倍2026年PE，市场对持续超预期已有较高期待。企业光通信需求超过产能，高密度连接器件采用加速，但当前收入仅反映横向扩展需求，纵向扩展和CPO采用仍可能带来额外上行空间。太阳能业务在晶圆转型期间仍实现90\%同比增长，管理层预计环比增长和利润率改善。
{\small\color{refgray}数据：企业光通信收入同比+65\%\par}
{\small\color{refgray}数据：Q3收入指引49-50亿美元，略低于共识\par}
{\small\color{refgray}数据：太阳能业务同比+90\%\par}
{\small\color{refgray}数据：估值超60倍2026年预期PE\par}
{\small\color{refgray}边际变化：市场关注点从AI光通信高增长转向Q3指引不及预期，估值从37倍降至31倍，反映对数据中心建设可持续性的担忧。\par}
\paragraph{Bernstein} 网络安全预算增速持续跑赢整体IT支出，但驱动因素已从合规与威胁响应转向AI。三位CISO一致认为，AI既是安全支出的最大催化剂，也迫使安全团队重新设计身份管理、运营模式和供应商策略。非人类身份（AI代理）管理成为焦点，安全运营正向自动化迁移，传统SIEM角色弱化。厂商整合集中在身份、端点、云安全等核心领域。
{\small\color{refgray}数据：保险行业CISO预算增长12-15\%\par}
{\small\color{refgray}数据：非人类身份管理成为首要关注\par}
{\small\color{refgray}数据：传统SIEM正被Agentic SOC模式替代\par}
{\small\color{refgray}数据：CrowdStrike、Wiz、Palo Alto、Okta被多次提及\par}
{\small\color{refgray}边际变化：AI从新兴趋势转变为重构安全项目的核心假设，身份安全市场从‘保护员工登录’转向‘管理机器与代理的权限生命周期’。\par}
\paragraph{GS} 印度非农就业中约13-15\%任务暴露于生成式AI自动化，但影响偏向互补而非替代。高暴露职业（常规文职、技术人员）面临替代风险，中等暴露职业（医疗、教育、专业服务）更可能被增强。替代不确定性集中在8-12\%就业，42-48\%就业可能被增强。AI可能使印度年劳动生产率增速提升0.1-0.8个百分点，基准情景0.4个百分点。
{\small\color{refgray}数据：13-15\%非农任务暴露于Gen-AI自动化\par}
{\small\color{refgray}数据：8-12\%就业面临替代不确定性\par}
{\small\color{refgray}数据：42-48\%就业可能被AI增强\par}
{\small\color{refgray}数据：劳动生产率增速提升0.1-0.8个百分点\par}
{\small\color{refgray}边际变化：从简单替代叙事转向‘增强多于替代’的结构性分化，服务业（医疗、教育、金融）成为AI增强主战场。\par}
\paragraph{MS} 中国AI计算从单芯片竞争转向系统级方案，规模互联技术成为核心载体。华为Atlas 950通过UnifiedBus 2.0连接1024颗NPU，设计目标8192颗；壁仞下一代BR2xx NPO架构支持1024颗GPU；中科曙光scaleX640单机架连接640颗加速器。芯片级限制使竞争焦点转向系统方案，光互联和服务器机架设计形成相对优势。
{\small\color{refgray}数据：华为Atlas 950连接1024颗NPU，目标8192颗\par}
{\small\color{refgray}数据：壁仞BR2xx NPO架构支持1024颗GPU\par}
{\small\color{refgray}数据：中科曙光scaleX640单机架640颗加速器\par}
{\small\color{refgray}数据：摩尔线程MTT C256单机架128颗GPU\par}
{\small\color{refgray}边际变化：WAIC 2026上芯片新品发布减少，超节点方案成为主角，竞争从芯片规格转向系统级互联能力。\par}
\paragraph{MS} 北美互联网板块7月24日当周加权下跌约7\%，标普500和纳斯达克分别下跌1\%和2\%。数字广告板块下跌7.8\%，共享宏观环境板块下跌8.0\%，旅游板块仅下跌2.8\%。META和GOOGL分别下跌约8\%，UBER和LYFT跌幅约9\%。AMZN、GOOGL、META当前对应2027年预期EPS的市盈率分别为20倍、21倍和18倍，比各自过去12个月滚动均值低12\%至25\%。
{\small\color{refgray}数据：互联网板块加权下跌7\%\par}
{\small\color{refgray}数据：数字广告板块下跌7.8\%\par}
{\small\color{refgray}数据：旅游板块仅下跌2.8\%\par}
{\small\color{refgray}数据：AMZN 2027年PE 20倍，低于均值25\%\par}
{\small\color{refgray}边际变化：市场在AI叙事、宏观预期与微观业绩三股力量交织下重新定价，旅游板块盈利可见度相对稳定，数字广告和共享经济板块面临更大不确定性。\par}
\paragraph{MS} 小米新电动汽车品牌SkyNomad将于7月30日在北京举办技术发布会，定位‘移动生活空间’，通过智能化重新定义车内空间。WLTC纯电续航最高380公里，馈电油耗最低5.7L/100km。截至2026年上半年已部署566辆测试车，累计路测超428万公里，覆盖中国300多个城市，测试温度范围-41°C至53°C，最高海拔5380米。
{\small\color{refgray}数据：WLTC纯电续航最高380公里\par}
{\small\color{refgray}数据：馈电油耗最低5.7L/100km\par}
{\small\color{refgray}数据：566辆测试车\par}
{\small\color{refgray}数据：累计路测超428万公里\par}
{\small\color{refgray}边际变化：小米汽车从研发阶段转向市场导入阶段，技术发布会成为评估其差异化能力的关键窗口。\par}
\paragraph{GS} ASMPT 2Q26营收49.36亿港元，环比+24\%，同比+52\%，超预期。毛利率42.4\%，环比+292bps，同比+284bps。运营利润率提升至26.5\%，运营利润7.86亿港元，环比+104\%。3Q26营收指引中值6.6亿美元，高于预期。热压键合（TCB）技术是AI基础设施周期的关键驱动力，已从领先IDM客户获得订单。
{\small\color{refgray}数据：2Q26营收49.36亿港元，同比+52\%\par}
{\small\color{refgray}数据：毛利率42.4\%，环比+292bps\par}
{\small\color{refgray}数据：运营利润率26.5\%\par}
{\small\color{refgray}数据：3Q26营收指引中值6.6亿美元\par}
{\small\color{refgray}边际变化：AI基础设施周期从概念验证进入实质订单阶段，TCB技术成为先进封装的关键驱动力。\par}
\paragraph{HSBC} MiniMax下一代文本模型M3 Pro预计2026年9-10月发布，参数规模约3万亿，仅约600亿参数被激活（激活率2\%）。使用8万亿数据训练已达到此前M3模型用40万亿数据训练的智能水平，训练效率提升约5倍。公司自2025年9月起自研算力集群，可运行超过1万芯片，开放平台毛利率目标为长期高双位数百分比。
{\small\color{refgray}数据：M3 Pro参数规模3万亿，激活率2\%\par}
{\small\color{refgray}数据：训练效率提升约5倍\par}
{\small\color{refgray}数据：自研集群可运行超1万芯片\par}
{\small\color{refgray}数据：日均处理数万亿推理token\par}
{\small\color{refgray}边际变化：从堆叠参数转向效率优化，稀疏激活和训练效率突破成为差异化优势，自建基础设施降低对云服务商依赖。\par}
\paragraph{GS} Resonac Holdings两个月内第二次宣布扩大HD媒体产能，追加约150亿日元投资，年产能再增2000万片，用于AI数据中心近线应用。公司首次明确披露已与硬盘制造商客户就研究额等额的涨价达成协议，在日本材料行业极为罕见。两轮扩产合计将使总产能从约1.6亿片提升至约2.3亿片，增幅约44\%。近线HDD出货容量2025-2030年复合增长率达23\%。
{\small\color{refgray}数据：追加150亿日元投资\par}
{\small\color{refgray}数据：年产能增加2000万片\par}
{\small\color{refgray}数据：总产能从1.6亿片增至2.3亿片，增幅44\%\par}
{\small\color{refgray}数据：近线HDD出货容量CAGR 23\%\par}
{\small\color{refgray}边际变化：从单纯扩产转向‘扩产+涨价’同步推进，验证独立供应商议价能力，盈利结构同步改善。\par}
\paragraph{JPM} 晶泰科技1H26盈利预警市场已有预期，去年同期有DoveTree 5100万美元首付款高基数。剔除BD收入后，非BD业务收入不低于2.5亿元，同比增长超65\%。AI for Science业务收入不低于1.8亿元，同比增长超120\%。研发费用不低于3.4亿元，同比增长超50\%，投入自动化实验室、智能体系统、多模态技术平台。
{\small\color{refgray}数据：非BD业务收入不低于2.5亿元，同比+65\%\par}
{\small\color{refgray}数据：AI for Science收入不低于1.8亿元，同比+120\%\par}
{\small\color{refgray}数据：研发费用不低于3.4亿元，同比+50\%\par}
{\small\color{refgray}数据：1H26营收3.8-4.0亿元\par}
{\small\color{refgray}边际变化：盈利预警反映研发投入主动加码而非基本面走弱，非BD服务形成自我维持增长曲线，AI平台化能力加速变现。\par}
\paragraph{MS} 2026年第二季度全球PC出货量约6620万台，环比持平，同比下降5\%，为连续九个季度增长后首次同比下滑。出货量下降但收入攀升，品牌厂商推动价格上涨速度快于需求下滑速度，形成‘量缩价升’局面。英特尔客户端计算业务收入同比增长15\%，驱动因素为更高平均售价而非出货量。服务器方面，英特尔数据中心业务收入同比增长59\%，需求远超供给，Q4仍将处于‘发货不足’状态。
{\small\color{refgray}数据：PC出货量6620万台，同比-5\%\par}
{\small\color{refgray}数据：英特尔客户端计算收入同比+15\%\par}
{\small\color{refgray}数据：英特尔数据中心收入同比+59\%\par}
{\small\color{refgray}数据：CPU/GPU部署比例接近1:1\par}
{\small\color{refgray}边际变化：PC市场从‘量价齐升’转向‘量缩价升’，高端化与涨价成为品牌厂商维持利润的主要手段；服务器供应瓶颈比预期更显著。\par}
\subsubsection*{关键数据}
\begin{itemize}
  \item 康宁企业光通信收入同比+65\%，但Q3指引略低于预期
  \item 印度非农就业中42-48\%可能被AI增强，仅8-12\%面临替代
  \item 华为Atlas 950连接1024颗NPU，设计目标8192颗
  \item 北美互联网板块一周下跌7\%，旅游板块仅跌2.8\%
  \item ASMPT 2Q26营收同比+52\%，毛利率42.4\%
  \item MiniMax M3 Pro训练效率提升约5倍，激活率仅2\%
  \item Resonac HD媒体总产能将增44\%，已与客户达成涨价协议
  \item 晶泰科技非BD业务收入同比+65\%，AI for Science收入同比+120\%
  \item 全球PC出货量同比-5\%，为九季度来首次下滑
  \item 英特尔数据中心收入同比+59\%，Q4仍将发货不足
\end{itemize}
\begin{figure}[htbp]
\centering
\includegraphics[width=0.88\linewidth]{figures/fig_020.jpg}
\caption{Exhibit 4 - Exhibit 4: AI supernode scale-up solutions Exhibit 5: Scale-up vs. scale-out networking \#\# China AI server scale-up development In China, electrical interconnect remains the preferred option within racks, while optical becomes mo}
\end{figure}
\begin{figure}[htbp]
\centering
\includegraphics[width=0.88\linewidth]{figures/fig_024.jpg}
\caption{Exhibit 4 - Exhibit 4: NTM EV/EBITDA for AMZN/GOOGL/META vs. Historical Averages Exhibit 5: Internet names fell -7\% last week (SPX/NDX -1/-2\%) Exhibit 6: NTM EV/EBITDA Multiples Are -14\%/-20\% vs. 5/10-Year Averages... Exhibit 7: ... and NT}
\end{figure}
\begin{figure}[htbp]
\centering
\includegraphics[width=0.88\linewidth]{figures/fig_112.jpg}
\caption{Exhibit 1 - Exhibit 1: Our preferred stocks are mostly related to AI and cloud, which we believe are more defensive vs. consumer electronics}
\end{figure}
\begin{figure}[htbp]
\centering
\includegraphics[width=0.88\linewidth]{figures/fig_039.jpg}
\caption{Exhibit 1 - Exhibit 1: Innovative design of SkyNomad to create a 'living space' inside the vehicle MS ASIA LIMITED+ Andy Meng, CFA}
\end{figure}
{\small\color{refgray}本节综合 11 篇研究；涉及 BofA、Bernstein、GS、MS、HSBC、JPM。}
\newpage
\subsection{半导体与硬件：存储定价分化与AI封装路线调整}
\textbf{存储芯片定价出现明显分化，DDR4及利基型NAND因结构性供需缺口在三季度迎来超预期涨价窗口，而主流DDR5涨幅趋缓；同时，AI加速器封装方案从4芯片向2芯片调整，EMIB-T技术提前商业化，产业链机会正在重新分配。}
\subsubsection*{主流共识}
\begin{itemize}
  \item AI基础设施需求持续扩张，数据中心冷却、光通信和低轨卫星等细分领域保持高增长。
  \item 存储芯片整体处于上行周期，三季度DRAM混合均价预计环比上涨13-18\%，NAND上涨10-15\%。
  \item 化合物半导体代工业务结构从手机芯片向卫星和AI数据中心转型，光学接收芯片量产进度超预期。
\end{itemize}
\subsubsection*{分歧与条件差异}
\begin{itemize}
  \item DDR4定价弹性分歧：MS认为DDR4三季度月涨幅可超20\%，而市场主流关注DDR5涨价放缓。
  \item GPU封装方案调整影响：NOM认为Rubin Ultra从4芯片改为2芯片将压低封装基板单价，但EMIB-T提前出货形成对冲。
  \item Skyworks苹果内容下滑幅度：MS报告显示下滑从20-25\%收窄至低十几个百分点，但长期影响仍需观察。
\end{itemize}
\subsubsection*{机构观点}
\paragraph{BofA} 希捷科技定价改善主要来自供需缺口扩大，而非单一客户优惠到期。增量毛利率从Q4的83.2\%升至Q1指引隐含的约88\%，近线硬盘产能已分配至2028年，客户规划延伸至2029年。上调FY27 EPS至32.06美元，FY28至54.06美元。
{\small\color{refgray}数据：FY26 Q4毛利率52.7\%，高于预期50.2\%\par}
{\small\color{refgray}数据：FY27营收预测174.7亿美元\par}
{\small\color{refgray}数据：增量毛利率约88\%\par}
{\small\color{refgray}数据：FY27 EPS从26.88上调至32.06美元\par}
{\small\color{refgray}边际变化：定价权从客户要求折扣转向供应商对增量产能拥有定价权，订单覆盖从季度延伸至FY27末。\par}
\paragraph{MS} Skyworks季度营收9.35亿美元超预期，苹果内容下滑幅度从20-25\%收窄至低十几个百分点，数据中心业务同比增速超50\%。但公司取消股息支付，毛利率受手机占比上升和金价上涨影响承压。Qorvo收购预计年内完成，长期毛利率目标50-55\%。
{\small\color{refgray}数据：营收9.35亿美元，超预期9.02亿\par}
{\small\color{refgray}数据：毛利率44.9\%，同比降217bps\par}
{\small\color{refgray}数据：EPS 1.08美元，超预期1.05\par}
{\small\color{refgray}数据：苹果内容下滑收窄至低十几个百分点\par}
{\small\color{refgray}边际变化：苹果内容下滑幅度显著收窄，但股息取消可能引发短期波动。\par}
\paragraph{NOM} NVIDIA Rubin Ultra从4芯片改为2芯片配置，因超大中介层与HBM4布线技术难度大。EMIB-T封装基板出货时间提前一年至2028年3月期，无需中介层降低生产难度。GPU封装基板月出货金额5月达281亿日元，每平方米价格升至136万日元。
{\small\color{refgray}数据：GPU封装基板月出货金额5月281亿日元\par}
{\small\color{refgray}数据：每平方米价格5月136万日元\par}
{\small\color{refgray}数据：EMIB-T出货提前至2028年3月期\par}
{\small\color{refgray}数据：Rubin Ultra改为2芯片配置\par}
{\small\color{refgray}边际变化：封装方案从4芯片改为2芯片，EMIB-T提前一年出货，正负因素大致抵消。\par}
\paragraph{GS} Carrier Q2业绩超预期，数据中心订单同比增超300\%，全年指引从15亿上调至20亿美元。美洲渠道库存同比下降25\%，下半年需求预期走强。对三菱电机和松下控股具有正面含义，但原材料成本上升影响利润率。
{\small\color{refgray}数据：Carrier Q2销售额63.51亿美元，超预期60.21亿\par}
{\small\color{refgray}数据：总订单同比增长40\%\par}
{\small\color{refgray}数据：数据中心订单同比增超300\%\par}
{\small\color{refgray}数据：美洲渠道库存同比下降25\%\par}
{\small\color{refgray}边际变化：数据中心冷却需求加速扩张，渠道去库存接近尾声，补库需求即将释放。\par}
\paragraph{JPM} 化合物半导体代工厂营收增长来自光学接收芯片量产提前和低轨卫星需求上升。产能利用率环比提升5ppt至65\%。光学业务2026/27年分别增长50\%以上和25\%以上，AI光通信收入占比2027年达两位数。上调2026年营收增长预测至30\%。
{\small\color{refgray}数据：产能利用率Q2环比升5ppt至65\%\par}
{\small\color{refgray}数据：Q3营收预计环比增长11\%\par}
{\small\color{refgray}数据：光学业务2026年增长50\%以上\par}
{\small\color{refgray}数据：AI光通信收入2026下半年达高个位数占比\par}
{\small\color{refgray}边际变化：光学业务从单一项目进入多项目并行阶段，AI光通信收入占比快速提升。\par}
\paragraph{MS} DDR4和利基型NAND定价弹性超预期。DDR4三季度月涨幅可超20\%，季度环比至少30\%，云计算厂商愿意签署1-2年长期协议。SLC/MLC NAND三季度环比涨50-60\%，NOR Flash涨30-40\%。主流DDR5涨幅趋缓，四季度预计收窄至3-8\%。
{\small\color{refgray}数据：DDR4三季度月涨幅超20\%\par}
{\small\color{refgray}数据：SLC/MLC NAND三季度环比涨50-60\%\par}
{\small\color{refgray}数据：NOR Flash三季度环比涨30-40\%\par}
{\small\color{refgray}数据：主流DRAM三季度均价环比涨13-18\%\par}
{\small\color{refgray}边际变化：DDR4定价预测从旧预测3-8\%上调至15-20\%，云计算厂商开始签署量价分离的长期协议。\par}
\subsubsection*{关键数据}
\begin{itemize}
  \item 希捷增量毛利率约88\%，FY27 EPS上调至32.06美元
  \item Skyworks苹果内容下滑从20-25\%收窄至低十几个百分点
  \item GPU封装基板月出货金额5月281亿日元，每平方米价格136万日元
  \item Carrier数据中心订单同比增超300\%，渠道库存降25\%
  \item 化合物半导体代工厂产能利用率Q2升至65\%，光学业务2026年增长50\%以上
  \item DDR4三季度月涨幅超20\%，云计算厂商签署1-2年长期协议
  \item SLC/MLC NAND三季度环比涨50-60\%，NOR Flash涨30-40\%
\end{itemize}
\begin{figure}[htbp]
\centering
\includegraphics[width=0.88\linewidth]{figures/fig_107.jpg}
\caption{Exhibit 1 - Exhibit 1: Latest pricing forecast for DRAM and NAND Exhibit 2: NOR flash demand and supply growth rates Exhibit 3: NOR flash demand growth and supply growth by}
\end{figure}
\begin{figure}[htbp]
\centering
\includegraphics[width=0.88\linewidth]{figures/fig_080.jpg}
\caption{Figure 1 - Figure 1: Win Semi annual revenue trend by application}
\end{figure}
\begin{figure}[htbp]
\centering
\includegraphics[width=0.88\linewidth]{figures/fig_033.jpg}
\caption{Exhibit 1 - \#\# OWNERSHIP POSITIONING Refinitiv; MSPB Content. Includes certain hedge fund exposures held with MSPB. Information may be inconsistent with or may not reflect broader market trends. Long/Short Ratio = Long Exposure / Short exposure. Sector \% of Total Net Expo}
\end{figure}
\begin{figure}[htbp]
\centering
\includegraphics[width=0.88\linewidth]{figures/fig_081.jpg}
\caption{Figure 1 - Figure 1: Win Semi annual revenue trend by application Figure 2: Win Semi revenue mix by application \% Figure 3: Win Semi capex trend}
\end{figure}
{\small\color{refgray}本节综合 6 篇研究；涉及 BofA、MS、NOM、GS、JPM。}
\newpage
\subsection{消费与零售：量价分化与结构性复苏}
\textbf{全球消费与零售板块呈现显著分化：可口可乐等成熟消费品公司依靠销量驱动实现超预期增长，而中国高端消费（老铺黄金、LVMH亚太）仍受金价波动和需求节奏影响，但结构性增长空间犹存。上游原奶周期拐点信号增强，下游现制饮品竞争格局正在重塑。}
\subsubsection*{主流共识}
\begin{itemize}
  \item 可口可乐二季度销量增长5\%为17年最强，驱动有机收入超预期，多家投行上调全年指引。
  \item LVMH时装皮具部门在连续七个季度下滑后重回正增长，但增长质量仍待观察。
  \item 中国原奶供需周期接近拐点，优然牧业和现代牧业双双扭亏为盈。
\end{itemize}
\subsubsection*{分歧与条件差异}
\begin{itemize}
  \item 老铺黄金二季度净利润低于预期30\%，Citi和JPM均下调预测，但分歧在于：Citi认为折扣和产品组合变化是主因，JPM则强调金价波动和季节性扰动，并看好下半年新品驱动反弹。
  \item 瑞幸加盟商同店GMV下滑10-15\%，但ASP回升；毛利率恢复至35\%，但人工成本上升和回本周期拉长（3年+）压制净利率，市场对盈利修复可持续性存在分歧。
  \item 联合利华核心日化业务（HPCCo）销量增长超预期400个基点，但食品业务疲软，市场关注增长能否延续至下半年。
\end{itemize}
\subsubsection*{机构观点}
\paragraph{Bernstein} 可口可乐二季度有机销售额增长6\%，远超预期3.5\%，销量增长5\%是核心驱动力，而非单纯提价。公司上调全年有机增长指引至5\%。联合利华二季度有机销售增长5.8\%，远超预期4.1\%，核心日化业务（HPCCo）销量增长比市场预期高出400个基点，但食品业务仅增0.2\%。
{\small\color{refgray}数据：可口可乐Q2有机销售额增长6\% vs 预期3.5\%\par}
{\small\color{refgray}数据：可口可乐Q2单位箱量增长5\%\par}
{\small\color{refgray}数据：联合利华Q2有机销售增长5.8\% vs 预期4.1\%\par}
{\small\color{refgray}数据：联合利华HPCCo有机销售增长7.6\%\par}
{\small\color{refgray}边际变化：联合利华市场等待两年的销量加速终于出现，且动力来自剥离冰淇淋后留下的核心日化业务。\par}
\paragraph{Citi} 老铺黄金二季度净利润比预期低约30\%，受收入和利润率双重挤压。折扣产品占比上升、促销力度加大、线上销售占比下降及运营去杠杆导致净利率从Q1的21.8\%降至18-20\%。下调2026-2028年净利润预测17-22\%。
{\small\color{refgray}数据：老铺黄金Q2净利润6.6-7.1亿元 vs 预期约9-10亿元\par}
{\small\color{refgray}数据：Q2净利率18-20\% vs Q1的21.8\%\par}
{\small\color{refgray}数据：Q2收入约33-39.5亿元\par}
{\small\color{refgray}边际变化：利润率调整来自销售结构与定价策略共同作用，而非单一需求变化。\par}
\paragraph{NOM} LVMH二季度亚太区（不含日本）有机收入同比增长4\%，中国客户支出同比持平，但国内消费意愿仍处历史高位。核心时装皮具部门两年后重回正增长（+1\%）。中国高端消费仍有结构性增长空间。
{\small\color{refgray}数据：LVMH Q2亚太区（不含日本）有机收入同比+4\%\par}
{\small\color{refgray}数据：LVMH时装皮具部门Q2有机收入同比+1\%\par}
{\small\color{refgray}数据：中国客户支出同比持平\par}
{\small\color{refgray}边际变化：中国消费者并未离场，只是购买节奏和渠道在变化——更集中在促销节点、更倾向于国内购买。\par}
\paragraph{GS} 可口可乐二季度有机收入增长6\%，超预期，销量增长5\%为17年最强（除疫情期）。上调全年有机增长指引至5\%。优然牧业上半年扭亏为盈（净利润7.39-9.03亿元），原奶价格降幅收窄、饲料成本下降、淘汰牛价格回升是主因，验证上游供需周期拐点。
{\small\color{refgray}数据：可口可乐Q2单位箱量增长5\%\par}
{\small\color{refgray}数据：可口可乐全年有机增长指引上调至约5\%\par}
{\small\color{refgray}数据：优然牧业1H26净利润7.39-9.03亿元 vs 去年同期亏损2.97亿元\par}
{\small\color{refgray}数据：原奶价格3.05元/公斤 vs 2025年末3.02元/公斤\par}
{\small\color{refgray}边际变化：优然牧业与之前现代牧业的盈利预告形成呼应，指向行业层面的供需再平衡。\par}
\paragraph{JPM} 老铺黄金二季度净利润低于预期，同店销售下滑14-50\%，但金价波动和季节性扰动是主因，下半年新品加速和门店升级有望推动反弹。LVMH时装皮具部门Q2重回正增长（+1\%），迪奥品牌加速跑赢部门平均。瑞幸加盟商同店GMV下滑10-15\%，但ASP回升，门店毛利率恢复至35\%，但人工成本上升和回本周期拉长（3年+）压制净利率。
{\small\color{refgray}数据：老铺黄金Q2净利润5.1-7.6亿元 vs JPM预期约8亿元\par}
{\small\color{refgray}数据：老铺黄金Q2同店销售增速-50\%至-14\%\par}
{\small\color{refgray}数据：LVMH时装皮具部门Q2有机收入同比+1\%\par}
{\small\color{refgray}数据：瑞幸加盟商门店毛利率约35\%\par}
{\small\color{refgray}边际变化：老铺黄金自Q2已加速新品发布、调整赠品门槛、优化服务流程，这些动作在下半年才会逐步反映在收入端。\par}
\paragraph{MS} 可口可乐二季度单位销量增长5\%，显著高于市场预期的2.2\%和过去九个季度平均的1.1\%。有机销售增长7\%，高于市场预期的5.0\%。Fairlife业务预计长期为有机增长贡献超过100个基点。公司全年指引偏保守。
{\small\color{refgray}数据：可口可乐Q2单位箱量增长5\% vs 市场预期2.2\%\par}
{\small\color{refgray}数据：可口可乐Q2有机销售增长7\% vs 市场预期5.0\%\par}
{\small\color{refgray}数据：Fairlife Q2收入同比+18\%\par}
{\small\color{refgray}数据：Fairlife预计长期贡献有机增长80-120个基点\par}
{\small\color{refgray}边际变化：可口可乐的有机销售增速已与同类大型消费品公司拉开明显差距，且这一差距在可预见的未来可能持续。\par}
\subsubsection*{关键数据}
\begin{itemize}
  \item 可口可乐Q2单位箱量增长5\%，为17年最强（除疫情期）
  \item 联合利华核心日化业务（HPCCo）销量增长比市场预期高出400个基点
  \item 老铺黄金Q2净利润低于预期约30\%，净利率从21.8\%降至18-20\%
  \item LVMH时装皮具部门Q2有机收入同比+1\%，结束连续七个季度下滑
  \item 优然牧业1H26扭亏为盈，净利润7.39-9.03亿元
  \item 瑞幸加盟商门店毛利率恢复至约35\%，但回本周期拉长至3年+
  \item 原奶价格3.05元/公斤，降幅收窄
  \item Fairlife Q2收入同比+18\%，预计长期贡献有机增长80-120个基点
\end{itemize}
\begin{figure}[htbp]
\centering
\includegraphics[width=0.88\linewidth]{figures/fig_001.jpg}
\caption{EXHIBIT 9 - EXHIBIT 9: Financial Forecast for KO EXHIBIT 10: Regression Analysis: P/E Multiple vs Organic Sales Growth as of Jul 28, 2026 \#\# APPENDIX - FINANCIAL FORECASTS EXHIBIT 11: KO Income Statement}
\end{figure}
\begin{figure}[htbp]
\centering
\includegraphics[width=0.88\linewidth]{figures/fig_054.jpg}
\caption{Exhibit 1 - Exhibit 1: Historical KO and PEP NTM P/E (2001 to Present) \#\# HundredX Insights We supplement our preview with data from HundredX, a mission-based data and insights company that takes an innovative approach to monitoring consumer}
\end{figure}
\begin{figure}[htbp]
\centering
\includegraphics[width=0.88\linewidth]{figures/fig_104.jpg}
\caption{Exhibit 2 - Exhibit 2: KO YoY US Scanner Sales Growth Consistently Much Stronger Than PEP/KDP Exhibit 3: KO YoY US Scanner Sales Growth Ahead of Mega-Cap Peers Regional Unit Case Results: KO beat consensus unit case growth in every region}
\end{figure}
\begin{figure}[htbp]
\centering
\includegraphics[width=0.88\linewidth]{figures/fig_106.jpg}
\caption{Exhibit 7 - Exhibit 7: KO Income Statement}
\end{figure}
{\small\color{refgray}本节综合 10 篇研究；涉及 Bernstein、Citi、NOM、GS、JPM、MS。}
\newpage
\subsection{中国汽车与出行：分化加速，盈利质量成核心标尺}
\textbf{中国汽车行业进入关键分化期，竞争焦点从价格转向产品、技术与渠道的综合能力，盈利质量而非销量增长成为评估企业的核心标尺。海外市场仍是结构性增长来源，但政策不确定性上升；成本压力在三季度集中显现，将加速车企分化。}
\subsubsection*{主流共识}
\begin{itemize}
  \item 2026年下半年国内需求缺乏宏观反弹动力，消费者信心是核心约束，政策支持仅为边际改善。
  \item 新产品周期是当前最可信的增长路径，但市场已推出约600款新车型，80-20法则下多数难以达到预期。
  \item 海外市场是结构性增长来源，自主品牌海外毛利率比国内高5-10个百分点，但欧洲关税调整风险上升。
\end{itemize}
\subsubsection*{分歧与条件差异}
\begin{itemize}
  \item 成本压力：存储芯片涨价导致单车成本增加约4000元，但垂直整合能力强的车企（如自研芯片可降本超1万元）与依赖外采的车企在三季度利润表现将显著分化。
  \item 技术路线：部分新势力加速自研芯片和布局人形机器人/自动驾驶出行服务，短期带来资本开支和亏损压力，而传统车企更侧重现有产品组合优化。
  \item 市场策略：宁德时代主动转向份额扩张，通过渗透低端市场换取增量利润，但市场对其单位利润稀释的解读存在分歧——管理层认为属产品组合变动而非结构性侵蚀。
\end{itemize}
\subsubsection*{机构观点}
\paragraph{Bernstein} 比亚迪在日本推出首款轻自动车BEV Racco，补贴后入门价199.5万日元，定价较预期更谨慎，年销目标1万辆。日本轻自动车市场年销约170万辆，占全国汽车销量40\%，由铃木和大发主导（合计份额近70\%）。即便完成目标，比亚迪在日整体份额仅0.3\%，轻自动车市场约0.6\%，短期冲击有限，但2026年日本BEV渗透率可能迎来变化节点，Racco恰逢窗口期。
{\small\color{refgray}数据：Racco补贴后入门价199.5万日元\par}
{\small\color{refgray}数据：日本轻自动车市场年销约170万辆，占汽车总销量40\%\par}
{\small\color{refgray}数据：铃木和大发合计份额近70\%\par}
{\small\color{refgray}数据：比亚迪年销目标1万辆，对应轻自动车份额约0.6\%\par}
{\small\color{refgray}边际变化：比亚迪首次进入日本轻自动车这一最大细分市场，定价低于同类BEV约10\%以上，但销售网络有限，短期对格局影响有限。\par}
\paragraph{JPM} 中国汽车行业竞争焦点已从价格转向产品+技术+渠道的综合能力。2026年下半年国内需求仍处调整，约600款新车型上市但多数难达预期。成本压力从二季度开始显现，存储芯片涨价使单车成本增加约4000元，三季度将是首个无库存缓冲的利润观察窗口。海外市场仍是结构性增长点，某自主品牌海外利润占比达80-90\%，毛利率比国内高5-10个百分点。
{\small\color{refgray}数据：2026年上半年中国市场推出约600款新车型\par}
{\small\color{refgray}数据：存储芯片涨价导致单车成本增加约4000元\par}
{\small\color{refgray}数据：某自主品牌海外利润占比达80-90\%\par}
{\small\color{refgray}数据：海外毛利率比国内高5-10个百分点\par}
{\small\color{refgray}边际变化：成本压力从二季度开始显现，三季度成为利润不确定性关键窗口；海外市场虽为结构性增长点，但欧洲插电混动车型面临加征关税风险。\par}
\paragraph{MS} 宁德时代转向积极份额扩张策略，通过产能和成本优势渗透低端市场。单位利润变化非结构性侵蚀，而是产品组合变动所致，同类产品利润率保持稳定，储能产品利润率甚至改善。储能出货量大幅上修：2026年从169GWh上调至218GWh（+29.1\%），2027-2028年分别上调24.5\%和24.6\%，储能成为增量利润核心驱动力。
{\small\color{refgray}数据：2026年储能出货量预测从169GWh上调至218GWh（+29.1\%）\par}
{\small\color{refgray}数据：2027年储能出货量上调24.5\%\par}
{\small\color{refgray}数据：2028年储能出货量上调24.6\%\par}
{\small\color{refgray}数据：2026年储能出货量预计同比增长79.8\%\par}
{\small\color{refgray}边际变化：宁德时代增长叙事从电动车电池龙头转向电池平台公司，储能业务规模效应将重塑利润结构；市场关注点应从每瓦时利润转向份额扩张带来的增量利润。\par}
\subsubsection*{关键数据}
\begin{itemize}
  \item Racco补贴后入门价199.5万日元
  \item 日本轻自动车市场年销约170万辆，占汽车总销量40\%
  \item 2026年上半年中国市场推出约600款新车型
  \item 存储芯片涨价导致单车成本增加约4000元
  \item 某自主品牌海外利润占比达80-90\%
  \item 宁德时代2026年储能出货量预测上调29.1\%至218GWh
  \item 2026年储能出货量预计同比增长79.8\%
  \item 动力电池出货量增速预计为31.7\%
\end{itemize}
\begin{figure}[htbp]
\centering
\includegraphics[width=0.88\linewidth]{figures/fig_002.jpg}
\caption{Exhibit 2 - EXHIBIT 1: BYD'S Racco starts at JPY 1,995k after subsidies EXHIBIT 2: The Racco is offered in three trims - the 200, 300 Plus, and 300 Premium}
\end{figure}
\begin{figure}[htbp]
\centering
\includegraphics[width=0.88\linewidth]{figures/fig_003.jpg}
\caption{Exhibit 6 - EXHIBIT 6: While the overall Japanese automobile market continues a gradual decline, the share of Kei-cars has been steadily increasing over the long term, maintaining a relatively resilient volume EXHIBIT 7: In the current Kei-c}
\end{figure}
\begin{figure}[htbp]
\centering
\includegraphics[width=0.88\linewidth]{figures/fig_005.jpg}
\caption{Exhibit 8 - EXHIBIT 8: Suzuki's sales have remained steady at 500K-600K units, while Daihatsu has recently seen a decline due to the certification fraud issue}
\end{figure}
\begin{figure}[htbp]
\centering
\includegraphics[width=0.88\linewidth]{figures/fig_004.jpg}
\caption{Exhibit 6 - EXHIBIT 6: While the overall Japanese automobile market continues a gradual decline, the share of Kei-cars has been steadily increasing over the long term, maintaining a relatively resilient volume EXHIBIT 7: In the current Kei-c}
\end{figure}
{\small\color{refgray}本节综合 3 篇研究；涉及 Bernstein、JPM、MS。}
\newpage
\subsection{工业与能源：结构性增长与周期博弈}
\textbf{全球工业与能源板块呈现结构性增长与短期周期博弈并存的格局。电力设备受益于电气化长期趋势，但短期估值回调；工业气体电子业务加速，定价能力回升；燃气轮机订单高可见度，核能等待政策催化；煤炭进口套利窗口持续，国内煤价温和上行。}
\subsubsection*{主流共识}
\begin{itemize}
  \item 电力设备长期增长逻辑未变，电气化趋势（如欧盟目标）支撑需求扩张。
  \item 工业气体电子气业务受益于半导体供应链，结构性增长信号明确。
  \item 燃气轮机供应偏紧，未来五年订单可见度高。
\end{itemize}
\subsubsection*{分歧与条件差异}
\begin{itemize}
  \item 华明装备：市场担忧竞争加剧，但MS认为竞争对手扩产可能指向更大的市场总规模。
  \item 斗山能源：核能业务等待年底政策节点，SMR订单能否落地是关键变量。
  \item 煤炭：进口套利窗口持续开放，但国内煤价温和上涨，市场预期进口增量会抑制涨幅。
\end{itemize}
\subsubsection*{机构观点}
\paragraph{MS} 华明装备自2026年初高点回调约54\%，市场对竞争加剧、美国市场进展及一季度业绩不及预期过度担忧。竞争对手MR扩产可能暗示市场总规模超预期扩张，欧盟电气化目标（2040年电气化率从23\%提升至46\%）将创造增量空间。海外扩张核心驱动力是成本优势和供应链多元化，美国业务预计2026年Q4有阶段性进展。2025-29年电力设备出口收入年均增长约30\%，海外收入占比从34\%提升至54\%。2027年后盈利增长有望重新加速。当前报价对应2027年预期市盈率约19倍，低于全球同业平均的30倍。
{\small\color{refgray}数据：华明装备自高点回调约54\%\par}
{\small\color{refgray}数据：2027年预期市盈率约19倍，全球同业平均30倍\par}
{\small\color{refgray}数据：欧盟目标2040年电气化率达46\%（目前23\%）\par}
{\small\color{refgray}数据：2025-29年电力设备出口收入年均增长约30\%\par}
{\small\color{refgray}边际变化：市场对竞争格局的担忧过度，竞争对手扩产可能指向更大的市场总规模，而非份额争夺。\par}
\paragraph{GS} Air Liquide 2026年上半年业绩超预期，第二季度气体与服务可比增长3.3\%，高于共识2.9\%。电子气业务增速从一季度的2.9\%跃升至9.5\%，几乎是共识预期4.6\%的两倍。未来12个月研究机会储备增至48亿欧元，其中超过50\%来自电子领域（此前约40\%）。工业商业务定价从一季度的3.4\%提升至5.0\%，连续两个季度改善。上半年经常性营业利润28.45亿欧元，利润率20.9\%，同比提升110个基点，高于全年指引的100个基点。经营现金流33.72亿欧元，同比增长13\%，自由现金流超预期14.5\%。
{\small\color{refgray}数据：2Q26气体与服务可比增长3.3\%（共识2.9\%）\par}
{\small\color{refgray}数据：电子气可比增长9.5\%（共识4.6\%）\par}
{\small\color{refgray}数据：未来12个月研究机会储备48亿欧元（1Q26为45亿欧元）\par}
{\small\color{refgray}数据：工业商定价5.0\%（1Q26为3.4\%）\par}
{\small\color{refgray}边际变化：电子气业务增速远超预期，项目储备中电子占比从40\%升至50\%以上，收入可见度提升。工业商定价连续两个季度回升，下游客户对价格调整接受度提升。\par}
\paragraph{JPM} 斗山能源二季度营业利润3140亿韩元，高于JPM预测2830亿韩元和市场预期2790亿韩元，超预期主要来自斗山山猫（美国建筑设备需求韧性）。母公司营业利润970亿韩元，略低于预期但基本符合。上半年累计获得12台燃气轮机、6台蒸汽轮机和5份长期服务协议，管理层表示2030-31年交付窗口谈判可见度高，未来五年供给偏紧。全年订单目标13.3万亿韩元，完成率53\%。核能方面，韩国第12次电力基本计划预计年底公布，西南部大型项目可能带来增量；美国有望获得西屋项目长周期材料订单。SMR目标下半年约1万亿韩元。2026-28年营业利润预测上调1-7\%。
{\small\color{refgray}数据：2Q26营业利润3140亿韩元（JPM预测2830亿，市场2790亿）\par}
{\small\color{refgray}数据：上半年累计12台燃气轮机、6台蒸汽轮机、5份LTSA\par}
{\small\color{refgray}数据：全年订单目标13.3万亿韩元，完成率53\%\par}
{\small\color{refgray}数据：SMR年度目标约1万亿韩元\par}
{\small\color{refgray}边际变化：燃气轮机订单高可见度，供应偏紧格局持续至2030-31年。核能等待年底政策节点，SMR订单能否落地是下半年关键。\par}
\paragraph{MS} 6月中国动力煤进口环比增长38.3\%，其中印尼煤占比53.8\%，发货量环比激增34.9\%。7月印尼煤进口套利窗口持续开放，进口量可能维持高位。国内动力煤价格温和上涨，秦皇岛港5500大卡725元/吨（周涨0.4\%），CCI 5500指数827元/吨（周涨1.3\%）。焦煤价格稳定，柳林4号主焦煤坑口价850元/吨，澳洲峰景焦煤230美元/吨（周跌1.3\%）。进口量持续高位可能限制国内煤价上行空间。
{\small\color{refgray}数据：6月中国动力煤进口环比+38.3\%\par}
{\small\color{refgray}数据：印尼煤占比53.8\%，发货量环比+34.9\%\par}
{\small\color{refgray}数据：秦皇岛5500大卡725元/吨，周涨0.4\%\par}
{\small\color{refgray}数据：CCI 5500指数827元/吨，周涨1.3\%\par}
{\small\color{refgray}边际变化：7月进口套利窗口持续开放，进口量可能维持高位，对国内煤价上行空间形成压制。\par}
\subsubsection*{关键数据}
\begin{itemize}
  \item 华明装备2027年预期市盈率约19倍，全球同业平均30倍
  \item Air Liquide 2Q26电子气可比增长9.5\%（共识4.6\%）
  \item 斗山能源上半年累计12台燃气轮机订单，全年目标完成53\%
  \item 6月中国动力煤进口环比+38.3\%，印尼煤占比53.8\%
  \item 欧盟目标2040年电气化率达46\%（目前23\%）
  \item Air Liquide 1H26 OIR margin 20.9\%，同比+110bps
  \item 斗山能源2Q26营业利润3140亿韩元，超预期12\%
  \item 秦皇岛5500大卡动力煤725元/吨，周涨0.4\%
\end{itemize}
\begin{figure}[htbp]
\centering
\includegraphics[width=0.88\linewidth]{figures/fig_023.jpg}
\caption{Exhibit 4 - Exhibit 4: Huaming's one-year forward P/E band}
\end{figure}
\begin{figure}[htbp]
\centering
\includegraphics[width=0.88\linewidth]{figures/fig_102.jpg}
\caption{Exhibit 1 - Exhibit 1: Chart of the Week – Coal Import By Countries Breakdown Exhibit 2 : Data for the week}
\end{figure}
\begin{figure}[htbp]
\centering
\includegraphics[width=0.88\linewidth]{figures/fig_103.jpg}
\caption{Exhibit 1 - Exhibit 1: Chart of the Week – Coal Import By Countries Breakdown Exhibit 2 : Data for the week MS does and seeks to do business with companies covered in MS. As a result, investors should be aware that the firm may have a conf}
\end{figure}
{\small\color{refgray}本节综合 4 篇研究；涉及 MS、GS、JPM。}
\newpage
\subsection{金融与数字资产：代币化加速与信用市场结构性分化}
\textbf{代币化资产正从试点走向运营部署，预计2030年达5.5万亿美元，但转型是渐进的；同时，AI资本开支正重塑超大规模企业债-股定价关系，长期信用利差波动性上升，而亚洲银行业ROA呈现结构性分化。}
\subsubsection*{主流共识}
\begin{itemize}
  \item 代币化资产规模将大幅增长，主要由公开市场标的（美国公司债和国债）驱动。
  \item 超大规模企业AI相关资本开支依赖债务融资，导致长期信用利差波动性上升。
  \item 亚洲银行业净息差和ROA呈现明显市场间分化，中国持续下行，印度和印尼保持高位。
\end{itemize}
\subsubsection*{分歧与条件差异}
\begin{itemize}
  \item 代币化转型路径：混合模式（新旧系统并行）vs. 完全链上化，短期主导权存在分歧。
  \item 超大规模企业自由现金流改善方式：削减资本开支（利好债券）vs. 经营现金流增长（利好股票），对资产定价影响不同。
  \item Snap信用基本面：现金流转化率大幅提升（63\%）与监管/竞争不确定性并存，市场定价是否充分存在分歧。
\end{itemize}
\subsubsection*{机构观点}
\paragraph{Citi} 代币化资产正从试点转向运营部署，当前全球规模约170亿美元，基准情景下2030年将达5.5万亿美元，乐观情景8.2万亿美元。增长主要由公开市场标的驱动，特别是美国公司债和国债。DTCC、NYSE、Nasdaq等基础设施巨头将代币化嵌入核心系统，受监管的链上货币（稳定币预计2030年达1.9万亿美元）提供结算基础，主要司法管辖区监管清晰度改善。但转型是渐进的，混合模式将在短期内主导，效率提升前运营复杂性先增加。
{\small\color{refgray}数据：全球代币化资产当前规模约170亿美元\par}
{\small\color{refgray}数据：2030年基准情景5.5万亿美元，乐观情景8.2万亿美元\par}
{\small\color{refgray}数据：稳定币预计2030年达1.9万亿美元\par}
{\small\color{refgray}数据：DTCC计划2026年底开始为期三年的代币化服务试点\par}
{\small\color{refgray}边际变化：从实验性试点转向运营部署，基础设施巨头（DTCC、NYSE、Nasdaq）入局是关键转折点，而非加密原生企业主导。\par}
\paragraph{GS} 超大规模企业长期信用利差波动性显著上升，信用利差对权益回报的beta正在回升，与过去几年长期信用beta走低的趋势形成对比。AI资本开支的大规模举债融资是根本原因，长期端利差波动性明显高于短端。GS数据显示，长期信用利差走阔与远期自由现金流回报率持续下降同步发生，且大部分偏离发生在过去几个季度的新债发行浪潮中。随着多年期发行周期延续，AI相关发行人的信用曲线将进一步陡峭化。
{\small\color{refgray}数据：超大规模企业长期信用利差对权益beta回升，与2022-2023年beta走低趋势相反\par}
{\small\color{refgray}数据：长期端利差波动性显著高于短端\par}
{\small\color{refgray}数据：自由现金流回报率持续下降与利差走阔同步\par}
{\small\color{refgray}数据：大部分beta偏离发生在过去几个季度新债发行浪潮中\par}
{\small\color{refgray}边际变化：过去长期信用beta走低因宏观利率调整和回报追逐，现在AI资本开支周期使长期端波动率重新抬头，beta回升。\par}
\paragraph{JPM} 亚洲各市场银行ROA正经历结构性分化。中国银行业ROA从2017年0.87\%持续降至2024年0.69\%，2025年预计0.65\%；印度从0.46\%低点回升至1.41\%；印尼稳定在2.8\%左右。分化背后是净息差、信用成本和收入结构的不同走向。中国NIM从2019年2.27\%降至2024年1.57\%，2025年预计1.47\%；印度和印尼NIM分别为4.16\%和5.74\%。香港NIM从2020年1.46\%低点回升至2024年1.91\%。
{\small\color{refgray}数据：中国银行业ROA：2017年0.87\% → 2024年0.69\% → 2025E 0.65\%\par}
{\small\color{refgray}数据：印度银行业ROA：从0.46\%低点回升至1.41\%\par}
{\small\color{refgray}数据：印尼ROA稳定在2.8\%左右\par}
{\small\color{refgray}数据：中国NIM：2019年2.27\% → 2024年1.57\% → 2025E 1.47\%\par}
{\small\color{refgray}边际变化：中国NIM持续下行反映贷款利率调整与存款成本刚性并存的结构性压力，而印度和印尼受益于高政策利率和信贷定价能力。\par}
\paragraph{JPM} Snap信用基本面被两股相反力量拉扯：竞争、技术迭代和监管审查（尤其用户年轻化使监管影响更直接）构成压力，但自由现金流转化率从2023年22\%升至2025年63\%，账上长期保持约20亿美元以上现金，净杠杆率不到1倍。债券今年迄今利差走阔120-130个基点，而科技高收益指数仅走阔74个基点，市场已在定价不确定性，但折扣是否足够取决于风险兑现速度。
{\small\color{refgray}数据：自由现金流转化率：2023年22\% → 2025年63\%\par}
{\small\color{refgray}数据：近三年平均自由现金流2.3亿美元\par}
{\small\color{refgray}数据：账上现金20亿美元以上，净杠杆率不到1倍\par}
{\small\color{refgray}数据：债券利差今年走阔120-130个基点，科技高收益指数走阔74个基点\par}
{\small\color{refgray}边际变化：自由现金流转化率大幅提升（22\%→63\%），但监管从罚款扩展到年龄限制等更严格手段，且Snap用户更年轻、现金更少，影响比同行更直接。\par}
\subsubsection*{关键数据}
\begin{itemize}
  \item 全球代币化资产当前约170亿美元，2030年基准情景5.5万亿美元
  \item 稳定币预计2030年达1.9万亿美元
  \item 中国银行业ROA从2017年0.87\%降至2025E 0.65\%
  \item 印度银行业ROA从0.46\%回升至1.41\%
  \item 中国NIM从2019年2.27\%降至2025E 1.47\%
  \item Snap自由现金流转化率从2023年22\%升至2025年63\%
  \item Snap债券利差走阔120-130个基点，科技高收益指数走阔74个基点
  \item 超大规模企业长期信用利差对权益beta回升，与2022-2023年趋势相反
\end{itemize}
\begin{figure}[htbp]
\centering
\includegraphics[width=0.88\linewidth]{figures/fig_008.jpg}
\caption{Figure 1 - \#\# 3. Increasing regulatory clarity A gradual shift toward clearer and more coordinated regulatory frameworks is strengthening the legal foundation for institutional adoption. However, this progress is uneven and increasingly highlights a double-edged dynamic.}
\end{figure}
\begin{figure}[htbp]
\centering
\includegraphics[width=0.88\linewidth]{figures/fig_051.jpg}
\caption{Exhibit 2 - Exhibit 3: Hyperscaler credit has sold off as the aggregate free cash flow yield advantage has eroded}
\end{figure}
\begin{figure}[htbp]
\centering
\includegraphics[width=0.88\linewidth]{figures/fig_084.jpg}
\caption{Figure 2 - Figure 2: Snapchat User Interface The business is predominantly funded by advertising revenue supplemented by growing subscriptions and other direct revenue products. Today, \textbackslash{}\textasciitilde{}80-90\% of the company's revenue comes from advertisi}
\end{figure}
\begin{figure}[htbp]
\centering
\includegraphics[width=0.88\linewidth]{figures/fig_009.jpg}
\caption{Figure 2 - At the same time, investor behavior and distribution models are shifting. There is growing demand to bring assets closer to investors, expanding access to digital-native capital pools, corporate treasurers, and new wealth segments seeking diversification acros}
\end{figure}
{\small\color{refgray}本节综合 4 篇研究；涉及 Citi、GS、JPM。}
\newpage
\subsection{医疗与生物科技：依沃西单抗全球临床路径与融资节奏}
\textbf{依沃西单抗全球注册路径的时间窗口因事件率节奏变化而拉长，Summit的现金消耗速度要求其提前补充资金，但关键OS数据读出节点（1H27）仍是改变竞争格局的核心事件。}
\subsubsection*{主流共识}
\begin{itemize}
  \item HARMONi-3鳞癌队列PFS分析预计在2H26完成，同时伴随一次早期OS中期分析，但该OS分析α分配极小，仅作为支持性数据。
  \item 真正有统计效力的OS中期分析安排在1H27，届时中位随访时间预计超过20个月，与HARMONi-6的21个月中位随访相近。
  \item Summit 2Q26现金约6.9亿美元，季度运营支出升至1.52亿美元，新提交3.8亿美元ATM融资计划以补充资金。
\end{itemize}
\subsubsection*{分歧与条件差异}
\begin{itemize}
  \item 非鳞癌队列PFS分析与鳞癌队列第二次OS分析的时间先后顺序尚未明确，存在不确定性。
  \item FDA对HARMONi额外OS分析数据的审评是否需要额外时间或召开ODAC会议，由监管机构决定，PDUFA日期存在调整可能。
\end{itemize}
\subsubsection*{机构观点}
\paragraph{JPM} 依沃西单抗全球临床开发路径因事件率低于预期而拉长。HARMONi-3鳞癌队列PFS分析预计2H26完成，同步进行早期OS中期分析（α分配极小，仅观察趋势）。真正有统计效力的OS中期分析在1H27，中位随访约20个月。Summit 2Q26现金6.9亿美元，季度运营支出1.52亿美元，新提交3.8亿美元ATM融资计划以支持临床推进。非鳞癌队列PFS分析预计1H27，但与鳞癌二次OS分析的先后顺序未明确。
{\small\color{refgray}数据：HARMONi-3鳞癌队列PFS分析预计2H26完成\par}
{\small\color{refgray}数据：早期OS中期分析α分配极小，仅支持性数据\par}
{\small\color{refgray}数据：1H27进行第二次OS中期分析，中位随访>20个月\par}
{\small\color{refgray}数据：Summit 2Q26现金6.907亿美元\par}
{\small\color{refgray}边际变化：首次披露HARMONi-3两次OS分析的具体时间表，并明确早期OS分析的性质为支持性而非统计显著，同时Summit提前启动融资以应对加速的现金消耗。\par}
\subsubsection*{关键数据}
\begin{itemize}
  \item HARMONi-3鳞癌队列PFS分析预计2H26完成
  \item 1H27进行第二次OS中期分析，中位随访>20个月
  \item Summit 2Q26现金6.907亿美元
  \item 季度运营支出1.518亿美元
  \item 新提交3.8亿美元ATM融资计划
\end{itemize}
{\small\color{refgray}本节综合 1 篇研究；涉及 JPM。}
\newpage
\subsection{区域市场与主题：人口拐点、太空竞赛与结构性分化}
\textbf{全球市场正经历多重结构性拐点：人口总量见顶预期提前、日本太空产业在规模与成本间博弈、韩国股市去杠杆后依赖回购重建支撑、以及中国手术机器人海外商业化进入收入确认阶段。这些主题均指向长期趋势的边际变化，而非短期周期性波动。}
\subsubsection*{主流共识}
\begin{itemize}
  \item 全球生育率持续低于更替水平，人口总量拐点可能早于联合国此前预测的2084年，在2050年代前后到来。
  \item 日本太空产业具备精密组件制造优势，但受限于高精度、低产量的成本约束，增长潜力需外部资本支撑。
  \item 韩国KOSPI回调主要由外资机械性减仓和散户杠杆去化驱动，而非企业基本面恶化。
\end{itemize}
\subsubsection*{分歧与条件差异}
\begin{itemize}
  \item 人口下降对消费需求的影响时点：部分观点认为2030年前后显著，另一部分认为影响更远。
  \item 日本太空公司能否突破规模瓶颈：乐观者认为海外发射服务和组件出口可带动增长，悲观者认为成本下降速度远不及美国和中国。
  \item 韩国市场下一阶段支撑来源：NOM强调企业回购将成为结构性需求，但市场对回购规模和持续性存疑。
\end{itemize}
\subsubsection*{机构观点}
\paragraph{GS} 全球人口已过峰值出生年份（2012年），低生育率从富裕国家扩散至全球三分之二人口，包括墨西哥、印度和拉丁美洲。核心驱动是养育孩子的机会成本持续上升，而现有政策工具效果有限。人口结构变化对产业需求的影响将在2030年前后变得显著，依赖年轻人口增长的消费品类需求基础将比预期更早受到侵蚀。
{\small\color{refgray}数据：2012年为全球峰值出生年份\par}
{\small\color{refgray}数据：全球三分之二人口生活在生育率低于2的国家\par}
{\small\color{refgray}数据：拉丁美洲整体生育率约1.8\par}
{\small\color{refgray}数据：以1.5生育率计算，每一代人规模缩减近一半\par}
{\small\color{refgray}边际变化：市场对人口问题的关注点从发达国家老龄化转向全球性低生育率，人口总量见顶时间预期提前约30年。\par}
\paragraph{MS} 日本太空产业处于关键节点：全球火箭发射次数过去五年年复合增长约25\%，但日本年均仅5次。日本在精密组件制造领域有全球竞争力，但受限于高精度、低产量的成本约束，缺乏规模导致成本难以下降，成本高企又限制商业应用扩展。MS对行业给出与大盘持平观点，认为多数公司需外部资本支撑前期投入。
{\small\color{refgray}数据：全球火箭发射次数过去五年年复合增长率约25\%\par}
{\small\color{refgray}数据：日本2024-2025年每年仅发射5次\par}
{\small\color{refgray}数据：日本卫星年发射量约15颗\par}
{\small\color{refgray}数据：QPS Holdings目标价2700日元\par}
{\small\color{refgray}边际变化：日本太空产业从依赖国内发射转向借助海外服务实现增长，但行业整体尚未进入自我循环阶段。\par}
\paragraph{MS} SpaceX星舰第13次飞行测试整体评价A-，在载荷部署、在轨发动机重启和飞船再入回收三个关键环节达到或超出预期。首次部署20颗生产型星链V3卫星，验证了星舰实际投放生产型卫星的能力。助推器着陆阶段出现硬着陆，但这是V3助推器第二次着陆尝试，比第12次有明显进步。马斯克表示第14次飞行将尝试首次捕获飞船。
{\small\color{refgray}数据：33台发动机全部成功点火，零发动机故障\par}
{\small\color{refgray}数据：首次部署20颗生产型星链V3卫星\par}
{\small\color{refgray}数据：约39分钟时成功完成单台发动机重新点火\par}
{\small\color{refgray}数据：飞船在印度洋实现最软海上溅落，船体完整\par}
{\small\color{refgray}边际变化：星舰从测试阶段走向运营阶段的路径更清晰，第14次飞行可能首次实现飞船捕获，是迄今最明确的进展信号。\par}
\paragraph{NOM} 韩国KOSPI从6月22日9,115点高点回落至7月24日6,691点，跌幅约27\%，主要由外资机械性减仓、杠杆ETF扩张和机构资金配置触及上限三股力量叠加造成，而非企业基本面变化。外资年内累计净卖出约158万亿韩元，杠杆ETF回报率-53\%。下一阶段市场再定价将依赖企业回购，预计2026至2028年回购规模从116万亿韩元增至328万亿韩元，约90\%来自两家大型半导体公司。
{\small\color{refgray}数据：KOSPI从9,115点跌至6,691点，跌幅约27\%\par}
{\small\color{refgray}数据：外资年内累计净卖出约158万亿韩元\par}
{\small\color{refgray}数据：杠杆ETF回报率-53\%，远高于KOSPI同期-22\%\par}
{\small\color{refgray}数据：预计2028年回购规模达328万亿韩元\par}
{\small\color{refgray}边际变化：市场下跌驱动力从基本面担忧转向资金面去杠杆，企业回购成为新的结构性需求来源。\par}
\paragraph{GS} 中国手术机器人公司海外商业化从订单验证进入收入加速确认阶段。微创机器人1H26收入同比增长约200-230\%，海外收入增速超450\%，毛利率提升超15个百分点。Toumai全球累计商业订单约300台，海外超240台。精锋医疗累计装机约160台，海外约100台，累计手术量超2万例。海外市场已成为关键盈利和增长驱动力，国内招标量1H26同比小幅下滑至约30台，采购均价降至1330万元。
{\small\color{refgray}数据：微创机器人1H26收入同比增长约200-230\%\par}
{\small\color{refgray}数据：海外收入增速超450\%\par}
{\small\color{refgray}数据：Toumai全球累计订单约300台，海外超240台\par}
{\small\color{refgray}数据：精锋医疗累计手术量超2万例\par}
{\small\color{refgray}边际变化：海外商业化从订单积累转向收入确认，手术量增速成为可持续收入的锚点。\par}
\paragraph{GS} 长鑫存储于7月27日在科创板上市，募资86亿美元，上市首日较发行价上涨466\%，市值达4840亿美元，成为A股最大科技股IPO及中国权益市场市值第二大公司。历史数据显示，大型IPO上市后短期表现强劲：首日平均回报34\%，21个交易日平均回报31\%，A股1个月超额回报49\%。指数层面呈现V型走势，发行前一周沪深300平均下跌2\%，上市后一周反弹2\%。
{\small\color{refgray}数据：募资86亿美元，为A股最大科技股IPO\par}
{\small\color{refgray}数据：上市首日较发行价上涨466\%\par}
{\small\color{refgray}数据：市值达4840亿美元\par}
{\small\color{refgray}数据：历史大型IPO首日平均回报34\%\par}
{\small\color{refgray}边际变化：长鑫IPO规模创纪录，历史数据显示大型IPO后市场呈V型走势，同行业公司上市后平均上涨2-3\%。\par}
\paragraph{GS} 康宁2Q26核心销售额47.38亿美元，超预期；光通信部门销售额20.72亿美元，同比增32\%，利润4.38亿美元，同比增77\%。企业网络业务同比增65\%，AI数据中心需求接近弹性较高。但股价下跌约19\%，反映市场对AI相关公司短期情绪谨慎。日本光通信股估值差距拉大：藤仓26倍、住友电工19倍、古河电工22倍，而康宁41倍，未来存在重新定价空间。
{\small\color{refgray}数据：康宁2Q26核心销售额47.38亿美元，超预期\par}
{\small\color{refgray}数据：光通信部门利润同比增77\%\par}
{\small\color{refgray}数据：企业网络业务同比增65\%\par}
{\small\color{refgray}数据：日本光通信股远期PE 19-26倍 vs 康宁41倍\par}
{\small\color{refgray}边际变化：康宁业绩超预期但股价下跌，市场对AI相关公司情绪转谨慎；日本厂商估值折价扩大，市场情绪稳定后可能重新定价。\par}
\paragraph{GS} LandMark 2026年6月营收4.21亿新台币，同比增128\%，环比增4\%，2Q26营收12.21亿新台币，同比增121\%，环比增35\%。增长主要来自800G和1.6T硅光模块强劲需求、InP衬底供应改善及产能扩张。公司从Aixtron取得约8.37亿新台币设备，董事会批准30亿新台币设备采购计划。硅光收入占比从2025年70\%升至2026年预估87\%，营业利润率从2023年-27\%升至2026年预估41\%。
{\small\color{refgray}数据：6月营收4.21亿新台币，同比增128\%\par}
{\small\color{refgray}数据：2Q26营收12.21亿新台币，同比增121\%\par}
{\small\color{refgray}数据：硅光收入占比2026年预估87\%\par}
{\small\color{refgray}数据：营业利润率从2023年-27\%升至2026年预估41\%\par}
{\small\color{refgray}边际变化：硅光渗透率快速提升，产能扩张计划明确，盈利拐点已现。\par}
\paragraph{NOM} 美联储7月议息会议前夕，NOM降低韩国与台湾利率交易仓位，维持对香港和中国大陆的不同方向头寸。中期偏空框架下，市场对美联储讨论已从降息转向加息，地缘风险溢价仍在。短期交易机会方面，印度利率在油价70美元附近时风险回报不对称，已平仓空头头寸，若油价回到70-75美元可能重新入场。韩国利率偏多逻辑仍在：KOSPI走弱和韩元贬值降低激进加息概率。
{\small\color{refgray}数据：韩国5年期NDIRS利率4.06\%\par}
{\small\color{refgray}数据：台湾5年期NDIRS利率2.345\%\par}
{\small\color{refgray}数据：印度5年期NDOIS接近200日移动均线\par}
{\small\color{refgray}数据：韩国央行预计8月和10月各加息一次\par}
{\small\color{refgray}边际变化：市场对美联储预期从降息转向加息，亚洲利率中期偏空但短期存在不对称交易机会。\par}
\paragraph{Citi} 数字工程服务商EPAM和Globant二季度有机收入增长预计落在指引区间内，但可自由支配IT支出尚未释放，AI带来定价不确定性。Globant前50大客户贡献5.2\%收入增长，但小客户仍谨慎。EPAM受墨西哥客户因关税缩减支出和中东局势影响。Citi预计EPAM将小幅下调全年有机固定汇率收入增长指引至2-4\%，Globant下调至-0.7-1.2\%。
{\small\color{refgray}数据：Globant前50大客户收入增长5.2\%\par}
{\small\color{refgray}数据：EPAM 1Q有机增长3.7\%\par}
{\small\color{refgray}数据：EPAM固定价格合同占比从15.1\%升至21.7\%\par}
{\small\color{refgray}数据：Globant AI Pods年经常性收入同比增长约290\%\par}
{\small\color{refgray}边际变化：AI未能带来增长加速，反而体现为定价不确定性；市场关注下半年增速能否兑现。\par}
\paragraph{Bernstein} 开云集团1H26财报沟通释放战略转型信心信号。Gucci在连续12个季度负增长后于2Q26首次录得正增长，由全价销售驱动而非折扣渠道。管理层强调价格弹性呈指数级变化，部分单品主动下调定价。新创意总监Demna首个完整系列Primavera将于7月底进店。Saint Laurent和Bottega Veneta已恢复增长，Balenciaga和McQueen仍在调整中。上半年净关店84家，全年计划超100家。
{\small\color{refgray}数据：Gucci 2Q26首次正增长，结束连续12个季度下滑\par}
{\small\color{refgray}数据：上半年净关店84家\par}
{\small\color{refgray}数据：运营成本同比下降\par}
{\small\color{refgray}数据：珠宝线和眼镜线增长强劲\par}
{\small\color{refgray}边际变化：开云从依赖创意叙事拉升品牌溢价转向通过务实定价重建消费者连接，价格弹性策略效果超预期。\par}
\paragraph{MS} 市场可能低估法国企业附加税被长期延续的可能性。MS将LVMH和爱马仕2027-2029年有效税率设定在2025年附加税实施后水平，比市场一致预期高300-400bps。LVMH每股盈利下调约6\%，爱马仕下调约3.5\%。法国政治日程密集，2026年秋季预算季、2027年总统选举及可能提前的议会选举将使政策不确定性持续到2027年中。爱马仕74\%产品在法国制造，税负敞口最大。
{\small\color{refgray}数据：LVMH每股盈利下调约6\%\par}
{\small\color{refgray}数据：爱马仕每股盈利下调约3.5\%\par}
{\small\color{refgray}数据：爱马仕74\%产品在法国制造\par}
{\small\color{refgray}数据：LVMH 2025年附加税贡献接近7亿欧元\par}
{\small\color{refgray}边际变化：市场共识未包含附加税长期化因素，MS将远期税率假设锁定在2025年水平，反映政策不确定性上升。\par}
\subsubsection*{关键数据}
\begin{itemize}
  \item 2012年为全球峰值出生年份
  \item 全球三分之二人口生活在生育率低于2的国家
  \item 日本年均火箭发射仅5次，全球年复合增长25\%
  \item KOSPI从9,115点跌至6,691点，跌幅27\%
  \item 微创机器人海外收入增速超450\%
  \item 长鑫存储募资86亿美元，首日涨466\%
  \item 康宁光通信利润同比增77\%，股价跌19\%
  \item LandMark硅光收入占比从70\%升至87\%
\end{itemize}
\begin{figure}[htbp]
\centering
\includegraphics[width=0.88\linewidth]{figures/fig_059.jpg}
\caption{Exhibit 1 - Exhibit 1: Any long-term birth rate that averages under two would cause depopulation World population in billions under different TFR scenarios 2. Peak populations coming in earlier than expected in both developed and developing}
\end{figure}
\begin{figure}[htbp]
\centering
\includegraphics[width=0.88\linewidth]{figures/fig_034.jpg}
\caption{Exhibit 4 - Exhibit 4: Global rocket launches Exhibit 5: Global satellite launches \#\# Japan's space industry is also growing}
\end{figure}
\begin{figure}[htbp]
\centering
\includegraphics[width=0.88\linewidth]{figures/fig_041.jpg}
\caption{Exhibit 1 - Exhibit 1: 1H26 bidding volume slightly below 1H25, due to slower hospital capex spending Exhibit 2: Procurement ASP declined to Rmb13.3mn in 1H26 as domestic brands gained market share Exhibit 3: Only 6 da Vinci systems were n}
\end{figure}
\begin{figure}[htbp]
\centering
\includegraphics[width=0.88\linewidth]{figures/fig_046.jpg}
\caption{Exhibit 1 - Exhibit 1: Top-50 Chinese IPOs across the 3 major listing locations in the history of Chinese equities \#\# 2) How did large IPOs trade in the past? Using the top-50 largest IPOs in China onshore and offshore markets (22 A shares,}
\end{figure}
{\small\color{refgray}本节综合 12 篇研究；涉及 Bernstein、MS、NOM、Citi、GS。}
\newpage
\subsection{中国结构性分化：工业利润、地产与电力需求}
\textbf{中国上半年经济呈现明显的结构性分化：工业利润高度集中于AI与上游原材料，下游消费持续承压；房地产市场新房与二手房成交分化加剧，一线城市新房回暖但二线波动大；电力需求增速受气温扰动回落，但高技术制造与充电服务用电增长强劲。}
\subsubsection*{主流共识}
\begin{itemize}
  \item 工业利润增长集中在技术和上游原材料行业，下游消费行业利润持续调整。
  \item 房地产市场新房与二手房成交分化，二手房市场相对稳定。
  \item 电力需求增速受气温因素影响，但高技术制造和充电服务用电增长强劲。
\end{itemize}
\subsubsection*{分歧与条件差异}
\begin{itemize}
  \item 关于工业利润持续性：共识认为利润增速对AI和上游原材料价格敏感度高，但分歧在于下半年利润增速是否会因全球AI资本开支变化或大宗商品价格回调而快速收窄。
  \item 关于房地产复苏：部分观点认为一线城市新房成交回暖是市场底部信号，但另一些观点认为去化率大幅提升是短期现象，价格信号尚未趋势性改善。
\end{itemize}
\subsubsection*{机构观点}
\paragraph{JPM} 上半年中国规模以上工业企业利润同比增长18.7\%，但月度增速从5月的21.1\%回落至6月的15.1\%。利润增量高度集中于电子制造（+96.9\%）和有色金属（+99.4\%），两者贡献了约93\%的利润增量，而下游消费品行业利润持续下滑（家具-52.7\%，汽车-19.5\%）。产成品库存同比增速攀升至9.5\%，应收账款增速升至8.1\%，企业资金压力加大。预计下半年“双速”格局延续，利润增速对AI和上游原材料价格敏感度极高。
{\small\color{refgray}数据：1-6月工业企业利润同比+18.7\%，6月单月+15.1\%（5月+21.1\%）\par}
{\small\color{refgray}数据：电子制造利润同比+96.9\%，有色金属利润同比+99.4\%\par}
{\small\color{refgray}数据：家具利润同比-52.7\%，汽车利润同比-19.5\%\par}
{\small\color{refgray}数据：产成品库存同比+9.5\%，应收账款同比+8.1\%\par}
{\small\color{refgray}边际变化：利润增速从5月的21.1\%回落至6月的15.1\%，库存和应收账款增速攀升，企业资金压力加大。\par}
\paragraph{JPM} 陆金所2Q26新增贷款同比增长4.6\%，扭转1Q26的-14.8\%下滑；领先资产质量指标改善，C-M3环比下降20bp至1.0\%，30天逾期率下降30bp至5.8\%。当前市净率仅约0.1倍，市值相当于2025年末现金储备的37\%，估值未充分反映运营复苏。贷款利率IRR为18.6\%，低于同业平均的21.8\%，监管收紧风险相对较小。8月中期业绩可能成为盈利可见度催化剂。
{\small\color{refgray}数据：2Q26新增贷款同比+4.6\%（1Q26为-14.8\%）\par}
{\small\color{refgray}数据：C-M3环比-20bp至1.0\%，30天逾期率-30bp至5.8\%\par}
{\small\color{refgray}数据：市净率约0.1倍，市值/2025年末现金=37\%\par}
{\small\color{refgray}数据：2025年新贷款加权平均IRR 18.6\%（同业平均21.8\%）\par}
{\small\color{refgray}边际变化：新增贷款恢复同比增长，资产质量领先指标改善，估值仍处于历史低位。\par}
\paragraph{JPM} 联想集团AI服务器订单管道已超210亿美元，6月季度AI收入占ISG总收入30\%以上。预计FY27 ISG收入同比增长约60\%至305亿美元，营业利润率从FY26的0.4\%提升至5\%。PC业务OPM维持在6.8\%左右，服务器利润改善速度超过PC不确定性。GB300服务器已开始向新云客户交付，多元供应商基础和ODM+品牌双模式带来竞争优势。
{\small\color{refgray}数据：AI服务器订单管道超210亿美元\par}
{\small\color{refgray}数据：6月季度AI收入占ISG收入>30\%\par}
{\small\color{refgray}数据：预计FY27 ISG收入同比+60\%至305亿美元\par}
{\small\color{refgray}数据：ISG营业利润率从FY26的0.4\%提升至FY27的5\%\par}
{\small\color{refgray}边际变化：AI服务器收入占比突破30\%，订单管道持续增长，ISG利润率从0.4\%向5\%跃迁。\par}
\paragraph{JPM} 携程反垄断整改聚焦商业行为而非结构性干预，罚款5.18亿元一次性计入2Q26。市场可能将下半年表现误读为结构性变化，但JPM认为周期因素是主因：行业酒店RevPAR同比下降6\%，约8个百分点放缓来自宏观因素，仅4个百分点来自监管调整。酒店仍愿意为携程的增量需求和转化效率付费，整改改变的是变现形式而非盈利基础。
{\small\color{refgray}数据：罚款5.18亿元一次性计入2Q26\par}
{\small\color{refgray}数据：行业酒店RevPAR同比-6\%\par}
{\small\color{refgray}数据：二季度约8ppt放缓来自宏观，4ppt来自监管调整\par}
{\small\color{refgray}数据：预计3Q26/4Q26收入预期下调6\%/4\%\par}
{\small\color{refgray}边际变化：反垄断整改聚焦商业行为而非结构性干预，市场可能过度解读为结构性变化，实际周期因素为主。\par}
\paragraph{MS} 截至7月26日当周，50城新房周度备案成交同比下滑4\%，较前一周的+3\%转负，主要受二线城市拖累（同比-11\%，前一周+14\%）。10城二手房成交同比+5\%，延续稳定。一线城市新房成交同比+13\%，去化率从4\%跃升至52\%。二手房挂牌价同比降幅16.0\%，略有收窄。新房与二手房累计增速差距扩大至15个百分点，市场结构分化加剧。
{\small\color{refgray}数据：50城新房周度成交同比-4\%（前一周+3\%）\par}
{\small\color{refgray}数据：二线城市新房成交同比-11\%（前一周+14\%）\par}
{\small\color{refgray}数据：10城二手房成交同比+5\%\par}
{\small\color{refgray}数据：一线城市去化率从4\%升至52\%\par}
{\small\color{refgray}边际变化：新房成交同比转负，二线城市波动剧烈；二手房成交增速稳定，新房与二手房分化加剧。\par}
\paragraph{MS} 6月全社会用电量同比增长3.7\%，较5月的6.9\%明显回落，主要受居民用电同比下降3.1\%拖累（气温偏低）。第二产业用电增长4.7\%，第三产业增长5.6\%，均较5月放缓。但高技术制造和充电服务用电增长强劲：1H26充电桩服务用电同比+56.9\%，IT服务（含AIDC）用电同比+44.0\%。1H26风光发电占比18.9\%，较2025年全年16.7\%提升，但利用小时数下降113小时。
{\small\color{refgray}数据：6月全社会用电同比+3.7\%（5月+6.9\%）\par}
{\small\color{refgray}数据：居民用电同比-3.1\%（5月+7.5\%）\par}
{\small\color{refgray}数据：1H26充电桩服务用电同比+56.9\%\par}
{\small\color{refgray}数据：1H26 IT服务用电同比+44.0\%\par}
{\small\color{refgray}边际变化：6月用电增速因气温因素回落，但高技术制造和充电服务用电增长强劲，结构性支撑明确。\par}
\subsubsection*{关键数据}
\begin{itemize}
  \item 1-6月工业企业利润同比+18.7\%，6月单月+15.1\%
  \item 电子制造利润同比+96.9\%，有色金属利润同比+99.4\%
  \item 陆金所2Q26新增贷款同比+4.6\%，C-M3环比-20bp至1.0\%
  \item 联想AI服务器订单管道超210亿美元，6月AI收入占ISG>30\%
  \item 携程罚款5.18亿元，行业RevPAR同比-6\%
  \item 50城新房周度成交同比-4\%，二线城市同比-11\%
  \item 6月全社会用电同比+3.7\%，居民用电同比-3.1\%
  \item 1H26充电桩服务用电同比+56.9\%，IT服务用电同比+44.0\%
\end{itemize}
\begin{figure}[htbp]
\centering
\includegraphics[width=0.88\linewidth]{figures/fig_069.jpg}
\caption{Figure 1 - Figure 1: Lufax is trading at a 0.1x fwd P/B, close to the historical low... Figure 2: ...while its new loan growth improved to \$+4.6\textbackslash{}\%\$ y/y in 2Q26 and outstanding loan balance contraction moderated to \$-13.5\textbackslash{}\%\$ y/y Figure 3:}
\end{figure}
\begin{figure}[htbp]
\centering
\includegraphics[width=0.88\linewidth]{figures/fig_073.jpg}
\caption{Figure 1 - Figure 1: Lenovo PC shipments and YoY mn units, \%}
\end{figure}
\begin{figure}[htbp]
\centering
\includegraphics[width=0.88\linewidth]{figures/fig_096.jpg}
\caption{Exhibit 1 - Exhibit 1: Primary weekly unit sales and four-week moving average – Total 50 cities MS ASIA LIMITED+ Stephen Cheung, CFA}
\end{figure}
\begin{figure}[htbp]
\centering
\includegraphics[width=0.88\linewidth]{figures/fig_068.jpg}
\caption{Figure 1 - Lufax is a leading technology-enabled platform focused on retail and SME financing. We upgrade Lufax to Overweight from Neutral, as its 2Q26 update showed an emerging fundamental recovery, with new loan growth returning to positive territory and leading asset }
\end{figure}
{\small\color{refgray}本节综合 6 篇研究；涉及 JPM、MS。}
\newpage
\section{辅助信号｜战略咨询}
波士顿咨询2026年全球读者调查显示，AI长期信仰未动摇，但投资者对AI溢价的容忍度下降，要求看到可量化的基本面改善。市场定价与预期之间的裂痕扩大，企业需从叙事转向执行。
\subsection{AI定价透支，投资者要求基本面兑现}
\textbf{市场对AI的乐观情绪仍在，但56\%的读者认为当前定价过于乐观，净乐观度高达41\%。投资者不再为AI叙事买单，转而要求增长、利润率与自由现金流的实际改善。}
\begin{itemize}
  \item AI开发与监管被45\%读者列为前三重要因素，较2024年上升22个百分点，增幅最大。
  \item 56\%的读者认为市场对AI定价过于乐观，净乐观度41\%，远高于其他宏观因素。
  \item 约三分之二读者认为当前估值和AI预期可能阻碍未来回报，估值倍数压缩将成为TSR主要风险。
  \item 未来公司TSR将更依赖增长、利润率扩张和自由现金流生成，而非估值扩张。
  \item 投资者对AI执行能力存疑，不愿为AI投入牺牲利润率，要求明确的ROI路径。
\end{itemize}
\newpage
\section{辅助信号｜智库/国际机构}
本板块整合了BIS、IMF、OECD、WTO等国际机构的最新研究，聚焦AI投资对宏观政策的冲击、加密资产执法框架的实操进展、以及全球发展融资与贸易规则的结构性变化。
\subsection{AI投资浪潮重塑宏观与政策环境}
\textbf{AI相关投资已达部分经济体GDP的1\%，且融资结构从内部现金流转向债务市场，增加了利率敏感性和金融透明度挑战；同时AI对总供给和总需求的双重冲击模糊了央行判断周期信号的依据。}
\begin{itemize}
  \item BIS报告显示，美国、澳大利亚数据中心及IT制造投资分别达GDP的0.8\%和超1\%，全球AI相关投资预计从5000亿美元增至2030年的3-4万亿美元，数据中心资本支出同期增长两倍。
  \item 融资结构显著变化：2025年AI企业通过债券市场筹集2430亿美元（2023年仅790亿美元），私人信贷对AI相关企业未偿贷款从2016年近零增至2025年2000亿美元，占私人信贷总额比重从不到2\%升至8\%。
  \item AI同时作用于总需求和总供给，模糊了央行赖以判断宏观环境状态的周期信号，增加了货币政策校准的难度。
  \item 贸易条件效应分化：高AI含量商品出口增速最快，但增长主要来自价格而非数量，且不同地区受益不均。
\end{itemize}
\begin{figure}[htbp]
\centering
\includegraphics[width=0.88\linewidth]{figures/fig_117.jpg}
\caption{Figure 1 - \#\# 1.1. Multidimensional data space of LLL To address these aims and purposes, we begin by delineating the conceptual space of LLL and its associated data environment. This provides the foundation upon which we construct the core framework guiding our examinat}
\end{figure}
\subsection{加密资产执法与发展融资的结构性瓶颈}
\textbf{南非9-12\%人口持有加密资产，IMF协助其建立跨部门执法框架，但技术差距仍是核心挑战；OECD报告显示私人融资动员创历史新高，但工具集中度高、低收入国家仅获8\%，结构性瓶颈未突破。}
\begin{itemize}
  \item 南非加密资产监管形成多机构协作机制（财政部、央行、金融监管局、税务局、金融情报中心），正推进OECD加密资产报告框架（CARF）实施。
  \item IMF研讨会揭示核心挑战是执法部门与犯罪者的技术差距，欺诈者使用混币器、跳转器和点对点交易所隐藏交易路径。
  \item 2024年官方发展融资干预撬动私人资本770亿美元，创2012年以来新高，但非洲和拉美合计吸收近六成资金，低收入国家仅获8\%。
  \item 担保工具贡献30\%、直接投资于公司占25\%、银团贷款占16\%，三种工具合计占比超七成，工具集中度偏高，市场波动时动员能力将承压。
\end{itemize}
\begin{figure}[htbp]
\centering
\includegraphics[width=0.88\linewidth]{figures/fig_127.jpg}
\caption{Figure 1 - \#\# Overview Bilateral providers MDBs Other multilaterals —Ref. ODA, DAC countries (right-hand axis) Note: Left-hand axis shows mobilised private finance by official development finance interventions in USD constant 2023 prices. Right-hand axis shows ODA by DAC}
\end{figure}
\begin{figure}[htbp]
\centering
\includegraphics[width=0.88\linewidth]{figures/fig_128.jpg}
\caption{Figure 1 - \#\# Three leveraging mechanisms accounted for the largest share of private finance mobilised for sustainable development Percentage of total \#\# Box 1.1. Leveraging mechanisms used to mobilise private finance for sustainable development Development finance provi}
\end{figure}
\subsection{贸易规则与新兴市场结构性改革压力}
\textbf{WTO裁定土耳其电动汽车附加关税违反GATT，区域贸易协定例外不适用于委内瑞拉；马来西亚增长放缓，财政整固、生产力提升与教育改革相互制约，改革窗口已开。}
\begin{itemize}
  \item WTO专家组认定土耳其对进口电动汽车及部分混合动力汽车征收的附加关税超出约束税率，违反GATT第II:1(b)条和第I:1条最惠国待遇原则。
  \item 土耳其对委内瑞拉车辆的关税豁免因未能满足GATT第XXIV条标准，构成对中国产品的歧视性待遇。
  \item 马来西亚2025年人均购买力平价GDP约3.88万美元，过去五年平均实际增长5.2\%，但公共债务占GDP比重已达70.1\%，化石燃料补贴规模较大，与碳定价努力矛盾。
  \item OECD建议马来西亚削减化石燃料补贴、重新引入增值税并配合定向转移支付，同时推进教育体系改革以支撑生产力目标。
\end{itemize}
\begin{figure}[htbp]
\centering
\includegraphics[width=0.88\linewidth]{figures/fig_122.jpg}
\caption{Figure 1 - By the 1980s and 1990s, Malaysia had become an important hub in global value chains, particularly in the electrical and electronics sectors. The economy demonstrated resilience during shocks such as the 1997–98 Asian Financial Crisis and the Global Financial C}
\end{figure}
\begin{figure}[htbp]
\centering
\includegraphics[width=0.88\linewidth]{figures/fig_124.jpg}
\caption{Figure 3 - StatLink https://stat.link/75ao6s \#\# Productivity growth has slowed Productivity growth has slowed (Figure 3) and, in recent years, has fallen short of Malaysia's own targets. Regulatory complexity and skills shortages continue to constrain firm performance an}
\end{figure}
\subsection{AI、技术与生产率}
\textbf{用于判断 AI 投资周期、生产率扩散和产业约束是否正在改变资产定价。}
\begin{itemize}
  \item 终身学习已成为应对人口老龄化、技术变革和技能需求演化的核心政策工具。但经合组织在一份新发布的报告中指出，当前国际数据系统对终身学习的测量仍停留在碎片化阶段。报告提出的核心判断是：现有测量体系主要围绕正式教育和特定年龄段的培训参与设计，无法捕捉学习在整个人生历程和多元场景中的真实分...
\end{itemize}
\begin{figure}[htbp]
\centering
\includegraphics[width=0.88\linewidth]{figures/fig_117.jpg}
\caption{Figure 1 - \#\# 1.1. Multidimensional data space of LLL To address these aims and purposes, we begin by delineating the conceptual space of LLL and its associated data environment. This provides the foundation upon which we construct the core framework guiding our examinat}
\end{figure}
\section{结语}
后续需验证的关键节点包括：法国预算审议结果（赤字是否突破5.9\%）、DDR4三季度涨价幅度是否兑现、中国汽车三季度成本压力对利润的实际影响、依沃西单抗非鳞癌队列PFS分析时间表、以及全球人口拐点对消费需求影响的实证数据。
\newpage
\section{覆盖概览}
投行/券商 64 篇；战略咨询 1 篇；智库/国际机构 6 篇。\par
投行覆盖：JPM 18篇 / MS 16篇 / GS 13篇 / NOM 6篇 / Bernstein 5篇 / Citi 3篇 / BofA 2篇 / HSBC 1篇\par
\section{Disclaimer}
{\scriptsize\color{refgray}
This community is a paid learning and information-sharing community created for general educational purposes only. The membership fee is charged solely for access to community discussions, educational content, organization, and operational costs. It is not an advisory fee, management fee, performance fee, brokerage fee, or compensation for personalized financial, investment, legal, tax, or professional advice.\par
I participate and share information solely as an individual and for educational discussion among community members. I am not acting as a financial adviser, investment adviser, broker, dealer, portfolio manager, legal adviser, tax adviser, accountant, fiduciary, or any other licensed professional adviser. Nothing shared in this community constitutes, and should not be interpreted as, financial advice, investment advice, legal advice, tax advice, a recommendation, solicitation, offer, endorsement, instruction, or guarantee to buy, sell, hold, short, trade, or invest in any security, cryptocurrency, fund, derivative, strategy, or other financial product.\par
All content shared in this community, including messages, discussions, links, files, charts, examples, market commentary, watchlists, AI-generated or AI-polished summaries, and third-party materials, is general, impersonal, and educational in nature. It does not take into account any individual's financial situation, investment objectives, risk tolerance, investment horizon, liquidity needs, tax status, legal restrictions, or suitability requirements. No content should be relied upon as suitable for any specific person, account, portfolio, or financial decision.\par
Information shared in this community may be incomplete, inaccurate, outdated, speculative, AI-generated, AI-edited, or based on third-party internet sources that have not been independently verified. I make no representation or warranty as to the accuracy, completeness, timeliness, reliability, or suitability of any information shared.\par
Financial markets involve substantial risk, including the possible loss of principal. Past performance, historical data, hypothetical examples, simulated results, backtests, personal opinions, or educational examples do not guarantee future results. Each member is solely responsible for conducting their own research, exercising independent judgment, and making their own decisions. Before making any financial, investment, legal, or tax decision, members should consult qualified and licensed professionals.\par
Participation in this community, payment of any membership fee, or communication with me or other members does not create any adviser-client, fiduciary, professional, contractual, agency, partnership, employment, or other advisory relationship. By joining or participating in this community, you acknowledge and agree that you are solely responsible for your own decisions, actions, transactions, profits, losses, risks, and consequences, and that I shall not be liable for any direct, indirect, incidental, consequential, financial, legal, tax, or other loss or damage arising from or related to any content, discussion, information, or material shared in this community.\par
}
\newpage
\begin{center}
{\Large 更多详情报告kcdesk.com}\par\vspace{0.8cm}
\includegraphics[width=0.58\linewidth]{\detokenize{../../prompts/zsxq_img.jpg}}
\end{center}
\end{document}
